% GENERATED FILE. Edit canonical lessons, then run npm run generate:books.
% canonical-source-hash: fnv1a64:09b63ad6a8000360
% canonical-lessons: GE-C269-jahrzehnt, GE-C269-jahrtausend, GE-C269-epoche, GE-C269-jahrelang, GE-C269-stundenlang, GE-R269-review-262-269

\chapter{Jahrzehnte}
\label{ch:ge-jahrzehnte}
\hlchaptermodality{269}

\begin{chapteropening}
\textbf{By the end of this chapter:} \emph{I can talk about long stretches of time: a decade, a millennium, an era, for years and for hours.}

Complete the last lesson of chapter 269: I can talk about long stretches of time: a decade, a millennium, an era, for years and for hours.
\end{chapteropening}

\section[das Jahrzehnt]{das Jahrzehnt — the decade}

\label{lesson:GE-C269-jahrzehnt}



\begin{quote}
Before the new one: say the German for the length, then the German for the deadline.
\end{quote}

\subsection*{You'll want to know: das Jahrzehnt}
\textbf{das Jahrzehnt} — \textquotedblleft{}the decade\textquotedblright{}.

\begin{quote}
Das Dorf hat sich in einem Jahrzehnt sehr verändert.
\end{quote}

\begin{quote}
\emph{The village changed a lot in a decade.}
\end{quote}

Plural: \emph{die Jahrzehnte}. Ten years, as a \emph{Jahrhundert} is a hundred.

\begin{cousinweb}
\emph{Jahrzehnt} joins \emph{Jahr} \textquotedblleft{}year\textquotedblright{} and \emph{-zehnt}, built on \emph{zehn} \textquotedblleft{}ten\textquotedblright{}; EtymDB lists Dutch \emph{jaartien} \textquotedblleft{}decade\textquotedblright{} as its cousin. Both parts are inherited Germanic words.
\end{cousinweb}

\begin{grammarlens}[title={What you've built}]
One more word for telling, explaining and describing at B1.
\end{grammarlens}

\subsection*{Your turn}
\begin{itemize}
\raggedright
  \item \emph{Say it:} \emph{das Jahrzehnt}
  \item \emph{Say it:} \emph{das Jahrzehnt}, once more
  \item \emph{Say it:} a sentence of your own with \emph{das Jahrzehnt}
  \item \emph{Recall:} say the German for the length, then the German for the deadline, then say \emph{das Jahrzehnt} again
\end{itemize}

\begin{tcolorbox}[breakable,colback=teal!4,colframe=teal!35!black,title={Before you move on}]
What does das Jahrzehnt mean? (\textquotedblleft{}The decade\textquotedblright{}.) Say it once more.
\end{tcolorbox}
\section[das Jahrtausend]{das Jahrtausend — the millennium}

\label{lesson:GE-C269-jahrtausend}



\begin{quote}
Before the new one: say the German for the deadline, then the German for the decade.
\end{quote}

\subsection*{You'll want to know: das Jahrtausend}
\textbf{das Jahrtausend} — \textquotedblleft{}the millennium\textquotedblright{}.

\begin{quote}
Die Kirche ist fast ein Jahrtausend alt.
\end{quote}

\begin{quote}
\emph{The church is almost a thousand years old.}
\end{quote}

Plural: \emph{die Jahrtausende}.

\begin{cousinweb}
\emph{Jahrtausend} joins \emph{Jahr} \textquotedblleft{}year\textquotedblright{} and \emph{tausend} \textquotedblleft{}thousand\textquotedblright{}, from Old High German \emph{dūsunt}, Proto-Germanic \emph{þūsundī} (EtymDB).
\end{cousinweb}

\begin{grammarlens}[title={What you've built}]
One more word for telling, explaining and describing at B1.
\end{grammarlens}

\subsection*{Your turn}
\begin{itemize}
\raggedright
  \item \emph{Say it:} \emph{das Jahrtausend}
  \item \emph{Say it:} \emph{das Jahrtausend}, once more
  \item \emph{Say it:} a sentence of your own with \emph{das Jahrtausend}
  \item \emph{Recall:} say the German for the deadline, then the German for the decade, then say \emph{das Jahrtausend} again
\end{itemize}

\begin{tcolorbox}[breakable,colback=teal!4,colframe=teal!35!black,title={Before you move on}]
What does das Jahrtausend mean? (\textquotedblleft{}The millennium\textquotedblright{}.) Say it once more.
\end{tcolorbox}
\section[die Epoche]{die Epoche — the era, the epoch}

\label{lesson:GE-C269-epoche}



\begin{quote}
Before the new one: say the German for the decade, then the German for the millennium.
\end{quote}

\subsection*{You'll want to know: die Epoche}
\textbf{die Epoche} — \textquotedblleft{}the era, the epoch\textquotedblright{}.

\begin{quote}
Das war eine schwierige Epoche für das Land.
\end{quote}

\begin{quote}
\emph{That was a difficult era for the country.}
\end{quote}

Plural: \emph{die Epochen}. Three syllables: \emph{E-po-che}.

\begin{cousinweb}
\emph{Epoche} was borrowed, through Late Latin \emph{epocha}, from Greek \emph{epokhē} \textquotedblleft{}a holding back, a fixed point in time\textquotedblright{}; EtymDB lists no parent, so this follows the usual account.
\end{cousinweb}

\begin{grammarlens}[title={What you've built}]
One more word for telling, explaining and describing at B1.
\end{grammarlens}

\subsection*{Your turn}
\begin{itemize}
\raggedright
  \item \emph{Say it:} \emph{die Epoche}
  \item \emph{Say it:} \emph{die Epoche}, once more
  \item \emph{Say it:} a sentence of your own with \emph{die Epoche}
  \item \emph{Recall:} say the German for the decade, then the German for the millennium, then say \emph{die Epoche} again
\end{itemize}

\begin{tcolorbox}[breakable,colback=teal!4,colframe=teal!35!black,title={Before you move on}]
What does die Epoche mean? (\textquotedblleft{}The era, the epoch\textquotedblright{}.) Say it once more.
\end{tcolorbox}
\section[jahrelang]{jahrelang — for years}

\label{lesson:GE-C269-jahrelang}



\begin{quote}
Before the new one: say the German for the millennium, then the German for the era.
\end{quote}

\subsection*{You'll want to know: jahrelang}
\textbf{jahrelang} — \textquotedblleft{}for years\textquotedblright{}.

\begin{quote}
Er arbeitete jahrelang in einer Fabrik.
\end{quote}

\begin{quote}
\emph{He worked in a factory for years.}
\end{quote}

Built the same way: \emph{tagelang}, \emph{wochenlang}, \emph{monatelang} — for days, weeks, months.

\begin{cousinweb}
\emph{jahrelang} joins \emph{Jahre} \textquotedblleft{}years\textquotedblright{} and \emph{lang} \textquotedblleft{}long\textquotedblright{}, both inherited Germanic words (EtymDB).
\end{cousinweb}

\begin{grammarlens}[title={What you've built}]
One more word for telling, explaining and describing at B1.
\end{grammarlens}

\subsection*{Your turn}
\begin{itemize}
\raggedright
  \item \emph{Say it:} \emph{jahrelang}
  \item \emph{Say it:} \emph{jahrelang}, once more
  \item \emph{Say it:} a sentence of your own with \emph{jahrelang}
  \item \emph{Recall:} say the German for the millennium, then the German for the era, then say \emph{jahrelang} again
\end{itemize}

\begin{tcolorbox}[breakable,colback=teal!4,colframe=teal!35!black,title={Before you move on}]
What does jahrelang mean? (\textquotedblleft{}For years\textquotedblright{}.) Say it once more.
\end{tcolorbox}
\section[stundenlang]{stundenlang — for hours}

\label{lesson:GE-C269-stundenlang}



\begin{quote}
Before the new one: say the German for the era, then the German for for years.
\end{quote}

\subsection*{You'll want to know: stundenlang}
\textbf{stundenlang} — \textquotedblleft{}for hours\textquotedblright{}.

\begin{quote}
Wir haben stundenlang auf den Bus gewartet.
\end{quote}

\begin{quote}
\emph{We waited for the bus for hours.}
\end{quote}

An adverb: it answers \emph{wie lange?}

\begin{cousinweb}
\emph{stundenlang} joins \emph{Stunden} \textquotedblleft{}hours\textquotedblright{} and \emph{lang} \textquotedblleft{}long\textquotedblright{}. \emph{Stunde} is inherited and has the old English cousin \emph{stound} \textquotedblleft{}a while\textquotedblright{} (EtymDB).
\end{cousinweb}

\begin{grammarlens}[title={What you've built}]
One more word for telling, explaining and describing at B1.
\end{grammarlens}

\subsection*{Your turn}
\begin{itemize}
\raggedright
  \item \emph{Say it:} \emph{stundenlang}
  \item \emph{Say it:} \emph{stundenlang}, once more
  \item \emph{Say it:} a sentence of your own with \emph{stundenlang}
  \item \emph{Recall:} say the German for the era, then the German for for years, then say \emph{stundenlang} again
\end{itemize}

\begin{tcolorbox}[breakable,colback=teal!4,colframe=teal!35!black,title={Before you move on}]
What does stundenlang mean? (\textquotedblleft{}For hours\textquotedblright{}.) Say it once more.
\end{tcolorbox}
\section[(dialogue)]{Ein erster Rückblick}

\label{lesson:GE-R269-review-262-269}



\begin{quote}
Say the German for at first and afterwards.
\end{quote}

\subsection*{Your turn}
Say these aloud:

\begin{itemize}
\raggedright
  \item at first, afterwards, as a result, last of all, the order
  \item after, before, as soon as, as long as, since then
  \item back then, once, the other day, recently, the day before yesterday
  \item meanwhile, by now, in the meantime, so far, long ago
  \item suddenly, unexpectedly, gradually, in time, at once
  \item to last, to pass, to go on, to continue, to interrupt
  \item the moment, the point in time, the period, the length, the deadline
  \item the decade, the millennium, the era, for years, for hours
\end{itemize}

\begin{tcolorbox}[breakable,colback=teal!4,colframe=teal!35!black,title={Before you move on}]
At first? (\textbf{zunächst}.) Afterwards? (\textbf{anschließend}.) As a result? (\textbf{daraufhin}.) Last of all? (\textbf{zuletzt}.) The order? (\textbf{die Reihenfolge}.) After? (\textbf{nachdem}.) Before? (\textbf{bevor}.) As soon as? (\textbf{sobald}.) As long as? (\textbf{solange}.) Since then? (\textbf{seitdem}.) Back then? (\textbf{damals}.) Once? (\textbf{einst}.) The other day? (\textbf{neulich}.) Recently? (\textbf{kürzlich}.) The day before yesterday? (\textbf{vorgestern}.) Meanwhile? (\textbf{inzwischen}.) By now? (\textbf{mittlerweile}.) In the meantime? (\textbf{währenddessen}.) So far? (\textbf{bisher}.) Long ago? (\textbf{längst}.) Suddenly? (\textbf{plötzlich}.) Unexpectedly? (\textbf{unerwartet}.) Gradually? (\textbf{allmählich}.) In time? (\textbf{rechtzeitig}.) At once? (\textbf{zugleich}.) To last? (\textbf{dauern}.) To pass? (\textbf{vergehen}.) To go on? (\textbf{andauern}.) To continue? (\textbf{fortsetzen}.) To interrupt? (\textbf{unterbrechen}.) The moment? (\textbf{der Augenblick}.) The point in time? (\textbf{der Zeitpunkt}.) The period? (\textbf{der Zeitraum}.) The length? (\textbf{die Dauer}.) The deadline? (\textbf{die Frist}.) The decade? (\textbf{das Jahrzehnt}.) The millennium? (\textbf{das Jahrtausend}.) The era? (\textbf{die Epoche}.) For years? (\textbf{jahrelang}.) For hours? (\textbf{stundenlang}.)
\end{tcolorbox}
